\chapter{Un metodo fondato sul campo}
\label{cap:metodo_kompatscher}

Il capitolo precedente si è chiuso con una domanda aperta: se non un percorso, cosa bisogna cercare? Una risposta concreta esiste, ed è stata sviluppata da due studiosi, Klaus Kompatscher e Nandi Maria Hrozny Kompatscher, che hanno lavorato per anni sul territorio alpino tra la valle dell'Adige e quella dell'Isarco, in Trentino e in Alto Adige, un'area diversa dal Veneto ma per molti versi confrontabile per ambiente, quota, e tipo di fauna cacciata \citep{kompatscher2007}.

Il loro lavoro merita un'attenzione particolare in questo libro, non solo per i risultati che offre, ma per il metodo con cui li ha ottenuti: un metodo raramente raccontato per intero nelle sintesi divulgative, che tendono a riportarne solo le conclusioni generali (i cacciatori sceglievano posizioni strategiche su creste e valichi) senza i numeri e i criteri che stanno dietro quella conclusione. Questo capitolo si basa sul testo integrale del loro articolo, pubblicato su \textit{Preistoria Alpina} nel 2007, comprensivo dell'appendice che elenca oltre trecento siti.

\section{Un lavoro nato dal camminare, non dal calcolare}

Il punto di partenza del metodo di Kompatscher e Hrozny Kompatscher non è un modello costruito a tavolino da una mappa altimetrica. È un lavoro di prospezione sistematica sul terreno, condotto a partire dal 1990, che affina un primo modello di mobilità già proposto da uno dei due autori oltre dieci anni prima \citep{kompatscher2007}. Anno dopo anno, gli autori hanno percorso a piedi ampie porzioni del territorio alpino, documentando ogni traccia di frequentazione preistorica trovata, con fotografie, schizzi dell'area con coordinate geografiche e orientamento del versante, e successivamente il posizionamento su cartografia IGM alla scala 1:25.000. Ogni sito individuato è stato inserito in un archivio strutturato, con la distanza dalla fonte d'acqua più vicina, l'ampiezza della visuale sul territorio circostante, la posizione (dominante o più raccolta) rispetto al paesaggio, la distanza da quello che gli autori hanno potuto ricostruire come il percorso principale di transito nella zona.

Gli stessi autori dichiarano esplicitamente, con un'onestà metodologica rara, il limite di fondo del loro approccio: procedono in ordine inverso rispetto alla tradizionale indagine archeologica, che si sviluppa dalla piccola alla grande scala (scavo, analisi, interpretazione, modello). Loro hanno iniziato dalla valutazione del territorio su larga scala, poi hanno condotto le prospezioni sistematiche, quindi individuato i siti, e solo alla fine sono giunti a una riflessione sul rapporto tra i siti e le aree circostanti. Il loro metodo, insomma, non è predittivo in senso stretto: è descrittivo, costruito osservando prima cosa esiste realmente sul terreno, e poi cercando i pattern statistici che ne emergono.

\section{Come si valuta un sito: due parametri territoriali}

Per la localizzazione dei siti su larga scala, gli autori hanno preso in considerazione due parametri, applicati sistematicamente a decine di casi.

Il primo è la valutazione topografica dell'ambiente circostante: l'area che i cacciatori mesolitici potevano sfruttare nell'arco di una giornata. Per calcolarla, il raggio d'azione è stato fissato in due ore di cammino dal campo, un valore che corrisponde a seimilaquattrocento metri su terreno pianeggiante (a un ritmo di ottocento metri ogni quarto d'ora) ma solo a ottocento metri su forti pendii (cento metri ogni quarto d'ora), con valori intermedi per pendenze intermedie. Le aree situate a quote inferiori di quattrocento metri rispetto al sito venivano escluse dal calcolo, perché considerate fuori dal raggio di sfruttamento pratico quotidiano. Questi parametri, va sottolineato, non sono stati assunti a priori, ma derivano da ripetute prospezioni e verifiche sul campo.

Il secondo parametro è la posizione strategica nel territorio, valutata contando il numero di direttrici che da un sito conducono verso territori di caccia adiacenti. La combinazione di questi due parametri, testata su una cinquantina di siti archeologicamente documentati, ha mostrato una buona concordanza con i calcoli teorici.

\subsection{Due casi che mostrano il limite del modello semplice}

Gli autori stessi offrono due esempi che mostrano quanto il criterio non sia meccanico. A Compatsch, sull'Alpe di Siusi, gli accampamenti dispongono di cinquantacinque chilometri quadrati di territorio circostante, di cui trentasei quasi pianeggianti: un'area vastissima secondo qualunque calcolo. Ma escursioni lunghe e promettenti erano possibili, da quei siti, solo in un'unica direzione, verso sud-est. In questo caso, un vasto raggio d'azione ha avuto un'importanza secondaria rispetto alla scarsità di direttrici: le zone ai margini occidentale e settentrionale dell'Alpe (il Puflatsch, lo Spitzbühel, il Monte Piz) non mostrano infatti insediamenti, nonostante l'ampio spazio teoricamente disponibile, perché i valori calcolati per le direttrici utili erano troppo bassi.

Al Passo Sella, al contrario, dove sono state individuate almeno trentotto testimonianze di frequentazione mesolitica, un gruppo di cacciatori poteva sfruttare un'area quotidiana di circa trentotto chilometri quadrati, di cui undici quasi pianeggianti, ma con la possibilità di raggiungere territori adiacenti senza superare grandi dislivelli lungo tre direttrici distinte. Questa condizione geomorfologica di luogo di passaggio obbligato trova conferma nell'alto numero di ritrovamenti, sia del Mesolitico antico sia di quello recente.

\section{Il tracciato del percorso e la sua verifica sul terreno}

Il tracciato del percorso ideale, secondo gli autori, si snoda prevalentemente nella fascia tra il limite del bosco e la prateria alpina, seguendo criteri di economia dello spostamento simili a quelli discussi nell'ottavo capitolo di questo libro: lo spostamento è più agevole sopra il limite del bosco che dentro una fascia di vegetazione fitta, un crinale con vegetazione aperta è più percorribile di una valle boschiva stretta, l'attraversamento di un versante è difficile se questo presenta vallecole trasversali, e un percorso centrale offre più possibilità di deviazione laterale verso altri territori di caccia.

Un esempio concreto viene dalle Alpi Sarentine, dove un percorso lungo un crinale centrale mostra una forte concentrazione di siti al Giogo dei Prati, a due terzi dell'itinerario, mentre i dossi laterali distanti dal tracciato principale risultano sistematicamente privi di ritrovamenti, con una sola eccezione degna di nota. Un secondo esempio, all'Alpe di Luson, mostra un percorso che procede da sella a sella, aggirando i dossi: le selle, per la loro posizione topografica, erano luoghi ideali per l'allestimento di bivacchi, perché permettevano di controllare un territorio ampio, mentre i dossi e i terrazzi più bassi risultano privi di ritrovamenti.

\section{Tre tipi di sito, distinti da cosa vi si trova}

Un contributo fondamentale del lavoro di Kompatscher e Hrozny Kompatscher è la distinzione, fondata su un criterio verificabile e non solo intuitivo, tra tre categorie di sito, basata sul confronto tra i manufatti in pietra ritrovati.

Il campo base è il luogo dove l'insieme litico indica una permanenza prolungata, associata ad attività diversificate: il segno diagnostico è un rapporto equilibrato tra strumenti e armature (le punte usate per le armi da lancio). Il sito di caccia o di avvistamento si trova sempre in posizione dominante rispetto alla visibilità sul territorio circostante, e mostra una netta prevalenza di armature, microbulini e microschegge, prodotte o riparate sul posto. I punti di sosta, infine, restituiscono pochi manufatti e si trovano spesso in posizioni con visuale ridotta: una categoria residuale, per siti che non rientrano nelle prime due.

Gli autori avvertono onestamente che il record archeologico non sempre permette di distinguere un utilizzo diacronico o diversificato di uno stesso sito: un luogo come il Passo Pampeago, discusso più avanti, potrebbe essere stato usato sia per brevi soste sia per permanenze più lunghe in momenti diversi.

\section{Le soglie quantitative sull'acqua e sulla visibilità}

Sull'acqua, il dato è particolarmente istruttivo, perché smentisce un'intuizione facile. Non è vero che i campi base si trovassero il più vicino possibile a una fonte d'acqua: gli autori trovano che la distanza tipica è di venti-ottanta metri, raramente inferiore. La ragione è pratica: il terreno immediatamente a ridosso di un corso d'acqua o di una sorgente tende a essere umido, poco stabile, meno adatto a un accampamento prolungato. I punti di semplice avvistamento, invece, potevano trovarsi molto più lontani dall'acqua, talvolta senza alcuna fonte nelle immediate vicinanze.

Sulla visibilità, gli autori partono da un presupposto concreto: un branco di ungulati alpini è riconoscibile a occhio nudo fino a una distanza di circa millecinquecento metri, e su questa base hanno costruito una scala a cinque livelli di qualità della visuale, da quasi assente a ottimale in tutte le direzioni, applicandola sistematicamente a trecentoventinove siti. Da questa analisi emerge un dato interessante che introduce una distinzione cronologica: nel Sauveterriano (il Mesolitico più antico) la buona visibilità aveva più importanza che nel Castelnoviano (la fase successiva), dove sembra che venisse data priorità ad altri criteri. Gli autori collegano questa differenza al continuo innalzamento del limite del bosco nel corso del Mesolitico, che rendeva più difficile osservare la selvaggina a grande distanza.

Un'analisi sull'esposizione del versante, condotta su ottantuno siti totali, mostra un pattern altrettanto interessante: per i quarantasette siti sauveterriani non è riconoscibile un'esposizione preferenziale, mentre per i trentaquattro siti del Mesolitico recente l'orientamento verso sud è nettamente dominante. Anche questo dato, secondo gli autori, potrebbe collegarsi al progressivo cambiamento ambientale che caratterizza il passaggio tra le due fasi.

\section{Cinque casi di studio esemplari}

Per illustrare come i criteri si combinino nella pratica, gli autori propongono cinque casi di studio dettagliati, scelti tra le aree a più alta densità di siti.

Il sito del Passo Pampeago rappresenta un caso raro in cui tutti i criteri sono soddisfatti in un unico punto: qui si incrociano tre percorsi verso vaste aree di caccia adiacenti, l'area sfruttabile misura quarantasei chilometri quadrati, l'approvvigionamento idrico è garantito da una sorgente che sgorga direttamente dalla sella, e il terreno quasi pianeggiante offre numerose possibilità per piantare un campo. Un piccolo spostamento in qualunque direzione farebbe crollare tutti questi valori contemporaneamente, il che spiega perché i versanti e le terrazze adiacenti, pur astrattamente idonei, non mostrino tracce di frequentazione. La concentrazione di manufatti su un'area di circa ottomila metri quadrati, attribuibili sia al Mesolitico antico sia a quello recente, conferma l'eccezionalità del luogo.

I siti sulla Cresta di Siusi mostrano invece la concentrazione lungo un percorso lineare: un sentiero corre in direzione est-ovest lungo un crinale allungato, con punti di discesa verso vasti territori di caccia su entrambi i lati. I campi base si collocano a media quota, nelle porzioni piane del versante, accettando sia una visibilità ridotta sia una distanza considerevole (settanta-ottanta metri) dalla sorgente d'acqua; i punti di avvistamento, invece, si trovano più in alto, in posizione dominante, con visibilità ottimale verso nord. Colpisce l'assoluta assenza di tracce archeologiche nelle aree orientate a sud, e la conferma che i cacciatori salivano a osservare la selvaggina solo nella misura strettamente necessaria: un dosso leggermente sopraelevato, pur vicino, risulta completamente privo di materiale litico.

Il Passo Sella mostra la logica dell'incrocio di percorsi: qui, a differenza di Pampeago dove tutto converge in un punto, esistono diverse possibilità di accampamento lungo i vari accessi al passo, con densità di siti che aumenta verso il centro dell'area, dove si ipotizza il bivio principale, e ritrovamenti isolati e sporadici nelle zone periferiche. I terreni situati più in alto rispetto al passo, nonostante ripetute prospezioni, sono rimasti privi di risultati.

I Laghi di Pennes mostrano l'influenza diretta della disponibilità d'acqua: il Passo Pennes stesso, pur trovandosi all'incrocio di tre percorsi importanti e dotato di ottima visibilità e terreno idoneo, ha restituito solo due ritrovamenti scarsi, quasi certamente per la mancanza di una fonte d'acqua sul passo. I gruppi umani si sono invece accampati a circa settecento metri di distanza, presso i laghetti, accettando la lontananza dal punto di incrocio pur di avere accesso all'acqua: tutti i siti si trovano lungo il bordo esterno del terrazzo, non direttamente sulla riva, un pattern che gli autori riconoscono anche in molti altri insediamenti simili.

Il rifugio Franz Senn, in Val di Obernberg, in Tirolo, mostra infine l'importanza di una posizione centrale rispetto a più valli: qui l'asse tra la valle principale e una valle laterale ha determinato la posizione dei siti, nonostante l'esistenza di terreni astrattamente favorevoli anche nei dintorni, dove però le ricerche non hanno restituito alcuna traccia.

\section{Un confronto con l'ambiente umano non montano}

Un aspetto del lavoro raramente riportato nelle sintesi è il confronto esplicito che gli autori fanno tra il territorio alpino e quello usato dai cacciatori-raccoglitori in ambienti più favorevoli, sulla base di studi etnografici. Studi su gruppi umani attuali che vivono di risorse sparse mostrano un raggio d'azione tipico di dieci-quindici chilometri per la caccia e cinque-dieci chilometri per la raccolta, valori superati solo in casi eccezionali perché un raggio maggiore riduce la resa della caccia e impedisce ai raccoglitori di tornare al campo in giornata.

Confrontando questi valori con la morfologia alpina, gli autori mostrano un divario netto: un raggio d'azione di quattro ore di cammino, che su terreno pianeggiante corrisponderebbe a dodici-sedici chilometri di distanza, in ambiente montano si traduce in un'area di ottanta-centocinquanta chilometri quadrati necessaria per ottenere lo stesso valore teorico di sfruttamento, appena il venti-trenta per cento dell'equivalente in pianura. Un esempio concreto, basato su un allineamento di sei comprensori di caccia (Passo Pampeago, Passo Sella, Cresta di Siusi e altri), mostra che la distanza tra il primo e il sesto punto, percorribile in circa trentatré ore di cammino su un territorio di seicentottantasei chilometri quadrati, richiederebbe solo sei ore e mezza in un ambiente pianeggiante equivalente: un fattore di circa cinque volte in più, dovuto quasi interamente alla morfologia del terreno, non alle assunzioni sul comportamento umano.

\section{Il dato più concreto: il percorso anulare}

L'elemento del lavoro di Kompatscher e Hrozny Kompatscher che più di ogni altro trasforma un quadro qualitativo in qualcosa di quasi tangibile è la ricostruzione di un ipotetico percorso anulare, motivato da un dato materiale preciso: la presenza di manufatti in selce sudalpina ritrovati oltre la catena principale delle Alpi, lontanissimi dagli affioramenti di origine, insieme a manufatti in cristallo di rocca i cui giacimenti si trovano lungo lo spartiacque alpino, rinvenuti sia nei siti dolomitici sia più a est. Gli autori ritengono improbabile che, dopo essersi procurati il cristallo di rocca, i cacciatori scendessero fino alla Val d'Adige per poi risalire nuovamente in quota: è più plausibile un percorso ad anello, con permanenza prolungata in montagna.

Un esempio concreto: partendo dalla conca di Bolzano con scorte di selce, un gruppo avrebbe potuto risalire verso nord fino all'Alpe Weitenberg, nelle Alpi Aurine, procurarsi lì il cristallo di rocca esaurita la selce, e fare ritorno passando più a est, attraverso le Dolomiti, lungo un percorso economico che tocca il Passo delle Erbe, il Passo Sella e la Cresta di Siusi. La durata di questo tragitto è stimata in circa novanta ore di cammino, che a un ritmo massimo di cinque-sei ore al giorno richiederebbero diciassette giorni complessivi. Ipotizzando una permanenza in quota di circa centocinquanta giorni all'anno, da giugno a fine ottobre, ne risulta una sosta media teorica di 8,3 giorni per ciascun sito lungo il percorso: un ritmo di vita quasi immaginabile giorno per giorno, che rende meno astratta l'idea stessa di un gruppo di cacciatori-raccoglitori in movimento stagionale prolungato.

\section{Cosa si può prendere in prestito, e cosa no}

Il metodo di Kompatscher e Hrozny Kompatscher non può essere trasferito di peso sul territorio veneto, per una ragione semplice quanto fondamentale: si basa su oltre trecento siti trovati con anni di prospezione sul campo, un lavoro che, come si è visto nella prima parte di questo libro, non è mai stato condotto con la stessa intensità sulle Alpi venete. Senza quei siti di partenza, non è possibile replicare il loro procedimento statistico, che parte proprio dall'osservazione di dove i siti reali si trovano.

Quello che si può prendere in prestito, invece, sono le soglie quantitative che il loro lavoro ha reso disponibili: la distanza tipica dall'acqua per un campo base, il raggio di visibilità utile per un punto di avvistamento, il modo di calcolare un'area di sfruttamento corretta per la pendenza del terreno, il confronto con i valori etnografici di riferimento, e la consapevolezza, altrettanto importante, che persino un'area vasta e teoricamente favorevole (come Compatsch) può risultare poco usata se manca la giusta combinazione di direttrici, mentre un'area più piccola ma ben collegata (come il Passo Sella) può concentrare decine di testimonianze. Questi elementi, nati dall'osservazione diretta in un contesto alpino comparabile, offrono un punto di riferimento concreto per orientare una ricerca sul territorio veneto, pur senza la pretesa di prevedere con certezza dove un sito si troverà.

Resta un'ultima domanda, a cui il prossimo capitolo dedica attenzione: con questi strumenti in mano, cosa è davvero ragionevole cercare sulle Alpi venete, e con quali limiti onesti bisogna presentare un simile lavoro di ricerca?
